\section{The quark--antiquark--gluon component of the outgoing state}
\label{sec:wf}

In this section we evaluate eq.~\eqref{eq:formal} explicitly. As in the massless
case~\cite{Iancu:2022gpw}, we first compute each contribution without scattering ($\hat S\to1$) in
transverse momentum space, we then Fourier transform to transverse coordinates, where the eikonal
scattering is diagonal, and we finally insert the Wilson lines. We treat in detail the emission by the
antiquark (figure~\ref{fig:qbar}) and deduce the emission by the quark (figure~\ref{fig:q}) from a
symmetry transformation.

\begin{figure}[t]
\centering
\begin{tikzpicture}[scale=0.92]
% (a)
\begin{scope}
\shock{0.55}
\draw[photon] (-2.3,0)--(-1.3,0);
\draw[quark] (-1.3,0)--(1.7,0.95) node[right]{\small $\x$};
\draw[quark] (-1.3,0)--(1.7,-0.95) node[right]{\small $\y$};
\draw[gluon] (-0.45,-0.27)--(1.7,-0.15) node[right]{\small $\z$};
\node at (-0.25,-0.62) {\small $\yp$};
\node at (-0.3,-1.6) {(a)};
\end{scope}
% (b)
\begin{scope}[xshift=5.1cm]
\shock{-0.15}
\draw[photon] (-2.3,0)--(-1.3,0);
\draw[quark] (-1.3,0)--(1.7,0.95) node[right]{\small $\x$};
\draw[quark] (-1.3,0)--(1.7,-0.95) node[right]{\small $\y$};
\draw[gluon] (0.55,-0.585)--(1.7,-0.15) node[right]{\small $\z$};
\node at (-0.75,-0.47) {\small $\yp$};
\node at (-0.3,-1.6) {(b)};
\end{scope}
% (c)
\begin{scope}[xshift=10.2cm]
\shock{0.55}
\draw[photon] (-2.3,0)--(-1.3,0);
\draw[quark] (-1.3,0)--(1.7,0.95) node[right]{\small $\x$};
\draw[quark] (-1.3,0)--(-1.3,-0.5);
\draw[line width=1.6pt] (-1.42,-0.25)--(-1.18,-0.25);
\draw[quark] (-1.3,-0.5)--(1.7,-0.95) node[right]{\small $\y$};
\draw[gluon] (-1.3,-0.5)--(1.7,-0.15) node[right]{\small $\z$};
\node at (-0.3,-1.6) {(c)};
\end{scope}
\end{tikzpicture}
\caption{Gluon emission by the antiquark. (a) The gluon is emitted before the collision with the
shock wave (grey band) and the three partons scatter. (b) The gluon is emitted after the collision and
only the quark--antiquark pair scatters, at the transverse positions $\x$ and $\yp$. (c) Instantaneous
graph: the photon decay and the gluon emission occur at the same light-cone time (the instantaneous
antiquark propagator is marked by a bar). The graphs in which the photon crosses the shock wave
before decaying are taken into account by the subtraction $\hat S\to\hat S-1$.}
\label{fig:qbar}
\end{figure}

\subsection{Gluon emission by the antiquark}
\label{sec:qbar}

\subsubsection{Emission before the collision}
\label{sec:before}

Inserting the matrix elements \eqref{eq:Hgamma2} and \eqref{eq:Hg}--\eqref{eq:Hg2} and the energy
denominators \eqref{eq:den1} and \eqref{eq:den3} in the first line of eq.~\eqref{eq:formal} with
$\hat S\to1$, we find
\begin{align}
\big|\gamma_T^i(Q,q^+,\qq)\big\rangle^{(a)}_{\qqg}=&-\frac{1}{(2\pi)^6}\int\dd^2\kt\,\dd^2\pt
\int_0^1\dd\zeta\int_0^{1-\zeta}\dd\xi\;
\frac{e\,e_f\,g\,q^+\,\zeta(1-\zeta)}{2\sqrt\xi\,\big(\kt^2+\ef^2\big)}
\nonumber\\
&\times\frac{\big[\varphi^{ij}_{\lambda_1\lambda}(\zeta)\,\kt^j+\mf\,\mu^i_{\lambda_1\lambda}\big]
\big[\tau^{mn}_{\lambda\lambda_2}(\xi,\eta)\,\pt^{\,n}+\xi\,\of\,\mu^m_{\lambda\lambda_2}\big]\,t^a_{\alpha\beta}}
{\zeta(1-\zeta)^2\big(\pt^{\,2}+\of^2\big)+\xi\eta\big(\kt^2+\ef^2\big)}
\nonumber\\
&\times\big|\bar q^\beta_{\lambda_2}(\eta q^+,\qq-\kk-\pp)\;g^a_m(\xi q^+,\pp)\;q^\alpha_{\lambda_1}(\zeta q^+,\kk)\big\rangle\,,
\label{eq:wf-mom}
\end{align}
where repeated helicity, polarisation and colour indices are summed over. The numerator is the product
of the two vertices in eq.~\eqref{eq:vertices} and contains four terms, which correspond to the four
helicity paths: no flip ($\varphi\,\tau$), flip at the gluon vertex ($\varphi\,\mu$), flip at the photon
vertex ($\mu\,\tau$), and double flip ($\mu\,\mu$).

To go to the transverse coordinate representation we write
\begin{equation}
\big|\bar q^\beta_{\lambda_2}(\qq-\kk-\pp)\,g^a_m(\pp)\,q^\alpha_{\lambda_1}(\kk)\big\rangle
=\int_{\x,\y,\z}e^{i\kk\cdot\x+i\pp\cdot\z+i(\qq-\kk-\pp)\cdot\y}\;
\big|\bar q^\beta_{\lambda_2}(\y)\,g^a_m(\z)\,q^\alpha_{\lambda_1}(\x)\big\rangle\,,
\end{equation}
and we express $\kk$ and $\pp$ in terms of $\kt$ and $\pt$. The phase becomes
\begin{equation}
e^{i\kk\cdot\x+i\pp\cdot\z+i(\qq-\kk-\pp)\cdot\y}=e^{i\w\cdot\qq}\;e^{i\R\cdot\kt}\;e^{-i\Y\cdot\pt}\,,
\label{eq:phase}
\end{equation}
where
\begin{equation}
\w=\zeta\x+(1-\zeta)\yp=\zeta\x+\eta\y+\xi\z\equiv\cc\,,\qquad
\yp\equiv\frac{\eta\y+\xi\z}{1-\zeta}\,,\qquad \R\equiv\x-\yp\,,\qquad\Y\equiv\y-\z\,.
\label{eq:coords}
\end{equation}
Here $\yp$ is the transverse position of the antiquark before the emission, which coincides with the
centre of energy of the final antiquark--gluon pair, $\cc$ is the centre of energy of the three-parton
system, which coincides with the position $\w$ of the photon, and $\R$ and $\Y$ are the transverse
separations of the daughter partons at the two vertices.

In the massless case the Fourier transform of eq.~\eqref{eq:wf-mom} is given by the well-known
formula~\cite{Beuf:2017bpd,Iancu:2022gpw}
\begin{equation}
\int\frac{\dd^2\kk}{2\pi}\frac{\dd^2\pp}{2\pi}\;
\frac{k^i\,p^j\;e^{i\kk\cdot\X+i\pp\cdot\bm{Z}}}{\big(\kk^2+\bar Q^2\big)\big[a(\kk^2+\bar Q^2)+\pp^2\big]}
=-\frac{X^iZ^j}{\bm{Z}^2}\;\frac{\bar Q\,K_1\big(\bar Q\sqrt{\X^2+a\bm{Z}^2}\big)}{\sqrt{\X^2+a\bm{Z}^2}}\,,
\label{eq:FT-massless}
\end{equation}
which involves a single modified Bessel function. With massive quarks, the two denominators of
eq.~\eqref{eq:wf-mom} contain two different scales, $\ef$ and $\of$, and each of the four helicity paths
has a different transform. We define four radial functions $\cG^{(\alpha\beta)}(\R,\Y)$, with
$\alpha,\beta\in\{0,1\}$ counting the powers of $\kt$ and $\pt$ in the numerator, by
\begin{align}
(1-\zeta)\int\frac{\dd^2\kt}{2\pi}\frac{\dd^2\pt}{2\pi}\;&
\frac{e^{i\kt\cdot\R}\,e^{-i\pt\cdot\Y}\;\big\{\kt^j\pt^{\,n},\;\kt^j,\;\pt^{\,n},\;1\big\}}
{\big(\kt^2+\ef^2\big)\Big[\zeta(1-\zeta)^2\big(\pt^{\,2}+\of^2\big)+\xi\eta\big(\kt^2+\ef^2\big)\Big]}
\nonumber\\
&=\Big\{R^jY^n\,\cG^{(11)},\;\;i\,R^j\,\cG^{(10)},\;\;-i\,Y^n\,\cG^{(01)},\;\;\cG^{(00)}\Big\}\,.
\label{eq:Gdef}
\end{align}
Using a Schwinger parametrisation of the two denominators and performing the Gaussian integrals (see
appendix~\ref{app:radial}), one of the two Schwinger integrals can be done in closed form and we obtain
\begin{equation}
\boxed{\;\cG^{(\alpha\beta)}(\R,\Y)=\frac{a^\beta}{2\zeta(1-\zeta)}\int_0^1\frac{\dd u}{u^{1+\beta}}
\left(\frac{A_u}{\Delta_u}\right)^{\frac{\alpha+\beta}{2}}K_{\alpha+\beta}\Big(\sqrt{A_u\,\Delta_u}\Big)\,,
\qquad\alpha,\beta\in\{0,1\}\,,\;}
\label{eq:Gab}
\end{equation}
with
\begin{equation}
A_u\equiv\ef^2+\frac{u\,\of^2}{a}=\zeta(1-\zeta)Q^2+\mf^2\left(1+\frac{u\,\zeta\,\xi}{\eta}\right)\,,\qquad
\Delta_u\equiv\R^2+\frac{a\,\Y^2}{u}\,,
\label{eq:ADelta}
\end{equation}
and $K_\nu$ the modified Bessel function of the second kind. Explicitly,
\begin{align}
\cG^{(11)}&=\frac{a}{2\zeta(1-\zeta)}\int_0^1\frac{\dd u}{u^2}\;\frac{A_u}{\Delta_u}\,
K_2\Big(\sqrt{A_u\Delta_u}\Big)\,,
\nonumber\\
\cG^{(10)}&=\frac{1}{2\zeta(1-\zeta)}\int_0^1\frac{\dd u}{u}\,\sqrt{\frac{A_u}{\Delta_u}}\;
K_1\Big(\sqrt{A_u\Delta_u}\Big)\,,
\nonumber\\
\cG^{(01)}&=\frac{a}{2\zeta(1-\zeta)}\int_0^1\frac{\dd u}{u^2}\,\sqrt{\frac{A_u}{\Delta_u}}\;
K_1\Big(\sqrt{A_u\Delta_u}\Big)\,,
\nonumber\\
\cG^{(00)}&=\frac{1}{2\zeta(1-\zeta)}\int_0^1\frac{\dd u}{u}\;K_0\Big(\sqrt{A_u\Delta_u}\Big)\,.
\label{eq:Gexplicit}
\end{align}
The integrals converge at $u\to0$ thanks to the exponential decay of the Bessel functions, since
$\Delta_u\to\infty$ there. The parameter $u$ has a simple interpretation: it is the fraction of the
total Schwinger ``time'' which is spent in the three-parton state. For $\mf=0$ one has $A_u=\bar
Q^2\equiv\zeta(1-\zeta)Q^2$, the $u$-integral in $\cG^{(11)}$ can be performed and
eq.~\eqref{eq:FT-massless} is recovered, see section~\ref{sec:massless}.

In terms of these functions, and after adding the Wilson lines for the three partons which cross the
shock wave and subtracting the no-scattering limit, the contribution of the graph in
figure~\ref{fig:qbar}(a) reads
\begin{align}
\big|\gamma_T^i(Q,q^+,\w)\big\rangle^{(a)}_{\qqg}=&-\frac{e\,e_f\,g\,q^+}{2(2\pi)^4}\int_{\x,\y,\z}
\int_0^1\dd\zeta\int_0^{1-\zeta}\dd\xi\;\frac{\zeta}{\sqrt\xi}\;\delta^{(2)}(\w-\cc)\;
\Psi^{im}_{\lambda_1\lambda_2}[\cG](\R,\Y)
\nonumber\\
&\times\Big[U^{ab}(\z)\,V(\x)\,t^b\,V^\dagger(\y)-t^a\Big]_{\alpha\beta}\;
\big|\bar q^\beta_{\lambda_2}(\eta q^+,\y)\;g^a_m(\xi q^+,\z)\;q^\alpha_{\lambda_1}(\zeta q^+,\x)\big\rangle\,,
\label{eq:wf-a}
\end{align}
where $V$ and $U$ are Wilson lines in the fundamental and in the adjoint representation,
\begin{equation}
V(\x)={\rm T}\exp\left\{ig\int\dd x^+\,t^a A^-_a(x^+,\x)\right\}\,,\qquad
U(\x)={\rm T}\exp\left\{ig\int\dd x^+\,T^a A^-_a(x^+,\x)\right\}\,,
\end{equation}
and where, for a generic set of four radial functions $F=\{F^{(11)},F^{(10)},F^{(01)},F^{(00)}\}$,
\begin{align}
\Psi^{im}_{\lambda_1\lambda_2}[F](\R,\Y)\equiv\sum_\lambda\Big[&
\varphi^{ij}_{\lambda_1\lambda}(\zeta)\,\tau^{mn}_{\lambda\lambda_2}(\xi,\eta)\,R^jY^n\,F^{(11)}
+i\,\xi\,\of\;\varphi^{ij}_{\lambda_1\lambda}(\zeta)\,\mu^m_{\lambda\lambda_2}\,R^j\,F^{(10)}
\nonumber\\
&-i\,\mf\;\mu^i_{\lambda_1\lambda}\,\tau^{mn}_{\lambda\lambda_2}(\xi,\eta)\,Y^n\,F^{(01)}
+\mf\,\xi\,\of\;\mu^i_{\lambda_1\lambda}\,\mu^m_{\lambda\lambda_2}\,F^{(00)}\Big]\,.
\label{eq:Psi}
\end{align}
The upper indices $i$ and $m$ are the polarisations of the photon and of the gluon. The four terms in
eq.~\eqref{eq:Psi} are the four helicity paths; the first one is the only one which survives when
$\mf\to0$.

The argument of the Bessel functions shows how the virtuality of the photon limits the transverse size
of the partonic fluctuation before the collision. In the massless case this size is the distance $\cD$
defined by
\begin{equation}
\cD^2(\x,\y,\z)\equiv\zeta(1-\zeta)\R^2+\frac{\xi\eta}{1-\zeta}\Y^2
=\zeta\eta(\x-\y)^2+\zeta\xi(\x-\z)^2+\eta\xi(\y-\z)^2\,,
\label{eq:D}
\end{equation}
with $Q\cD\lesssim1$. In the massive case one has $\sqrt{A_u\Delta_u}\geq\ef R$ for all $u$, with
$\ef>\mf$: the size of the quark--antiquark pair is limited by the Compton wavelength of the heavy
quark even in the photoproduction limit $Q^2\to0$, where all the functions \eqref{eq:Gexplicit} remain
finite.

\subsubsection{Emission after the collision}
\label{sec:after}

For the graph in figure~\ref{fig:qbar}(b) the only differences with respect to the previous case are
one energy denominator and the Wilson lines. According to eq.~\eqref{eq:formal}, the denominator
$E_{\qqg}-E_\gamma$ is replaced by $E_{\qqbar}-E_{\qqg}$, that is,
\begin{equation}
\frac{1}{\zeta(1-\zeta)^2\big(\pt^{\,2}+\of^2\big)+\xi\eta\big(\kt^2+\ef^2\big)}\;\longrightarrow\;
-\frac{1}{\zeta(1-\zeta)^2\big(\pt^{\,2}+\of^2\big)}\,.
\label{eq:replace}
\end{equation}
The two transverse integrations now factorise. Using
\begin{align}
\int\frac{\dd^2\kt}{2\pi}\;\frac{e^{i\kt\cdot\R}\,\{\kt^j,1\}}{\kt^2+\ef^2}
&=\left\{\frac{iR^j}{R}\,\ef K_1(\ef R)\,,\;K_0(\ef R)\right\}\,,
\nonumber\\
\int\frac{\dd^2\pt}{2\pi}\;\frac{e^{-i\pt\cdot\Y}\,\{\pt^{\,n},1\}}{\pt^{\,2}+\of^2}
&=\left\{-\frac{iY^n}{Y}\,\of K_1(\of Y)\,,\;K_0(\of Y)\right\}\,,
\label{eq:FT-fact}
\end{align}
with $R\equiv|\R|$ and $Y\equiv|\Y|$, we define the analogue of eq.~\eqref{eq:Gdef},
\begin{align}
(1-\zeta)\int\frac{\dd^2\kt}{2\pi}\frac{\dd^2\pt}{2\pi}\;&
\frac{e^{i\kt\cdot\R}\,e^{-i\pt\cdot\Y}\;\big\{\kt^j\pt^{\,n},\;\kt^j,\;\pt^{\,n},\;1\big\}}
{\big(\kt^2+\ef^2\big)\,\zeta(1-\zeta)^2\big(\pt^{\,2}+\of^2\big)}
\nonumber\\
&=\Big\{R^jY^n\,\cA^{(11)},\;\;i\,R^j\,\cA^{(10)},\;\;-i\,Y^n\,\cA^{(01)},\;\;\cA^{(00)}\Big\}\,,
\label{eq:Adef}
\end{align}
with the factorised functions
\begin{align}
\cA^{(11)}&=\frac{\ef K_1(\ef R)}{\zeta(1-\zeta)\,R}\;\frac{\of K_1(\of Y)}{Y}\,,
&\cA^{(10)}&=\frac{\ef K_1(\ef R)}{\zeta(1-\zeta)\,R}\;K_0(\of Y)\,,
\nonumber\\
\cA^{(01)}&=\frac{K_0(\ef R)}{\zeta(1-\zeta)}\;\frac{\of K_1(\of Y)}{Y}\,,
&\cA^{(00)}&=\frac{K_0(\ef R)}{\zeta(1-\zeta)}\;K_0(\of Y)\,.
\label{eq:Aexplicit}
\end{align}
Each of them is the product of a leading-order photon wave function, $\ef K_1(\ef R)$ or $K_0(\ef R)$,
and of an emission kernel off the massive antiquark, $\of K_1(\of Y)/Y$ or $K_0(\of Y)$. The massless
emission kernel $Y^n/Y^2$, which has a power-like tail at large $Y$, is thus replaced by
$Y^n\of K_1(\of Y)/Y$, which decays exponentially for $Y\gtrsim1/\of$: this is the dead-cone
effect in transverse coordinate space.

In this graph only the intermediate quark--antiquark pair, at transverse positions $\x$ and $\yp$,
scatters, and the colour matrix of the emission vertex stands to the right of the Wilson lines. We get
\begin{align}
\big|\gamma_T^i(Q,q^+,\w)\big\rangle^{(b)}_{\qqg}=&+\frac{e\,e_f\,g\,q^+}{2(2\pi)^4}\int_{\x,\y,\z}
\int_0^1\dd\zeta\int_0^{1-\zeta}\dd\xi\;\frac{\zeta}{\sqrt\xi}\;\delta^{(2)}(\w-\cc)\;
\Psi^{im}_{\lambda_1\lambda_2}[\cA](\R,\Y)
\nonumber\\
&\times\Big[V(\x)\,V^\dagger(\yp)\,t^a-t^a\Big]_{\alpha\beta}\;
\big|\bar q^\beta_{\lambda_2}(\eta q^+,\y)\;g^a_m(\xi q^+,\z)\;q^\alpha_{\lambda_1}(\zeta q^+,\x)\big\rangle\,.
\label{eq:wf-b}
\end{align}
The Wilson-line structure of eq.~\eqref{eq:wf-b} is obtained from that of eq.~\eqref{eq:wf-a} by
letting $\y,\z\to\yp$, because $U^{ab}(\yp)V(\x)t^bV^\dagger(\yp)=V(\x)V^\dagger(\yp)t^a$. In
contrast with the massless case, however, the radial functions $\cA^{(\alpha\beta)}$ are not simply the
coincidence limit of the functions $\cG^{(\alpha\beta)}$; they only share their short-distance
singularities, see eq.~\eqref{eq:UV} below.

\subsubsection{Instantaneous contribution}
\label{sec:inst}

For the graph in figure~\ref{fig:qbar}(c) there is a single energy denominator, $E_{\qqg}-E_\gamma$, and
the matrix element \eqref{eq:Hinst}. In momentum space and without scattering,
\begin{align}
\big|\gamma_T^i\big\rangle^{({\rm inst})}_{\qqg}=&\frac{1}{(2\pi)^6}\int\dd^2\kt\,\dd^2\pt
\int_0^1\dd\zeta\int_0^{1-\zeta}\dd\xi\;\frac{e\,e_f\,g\,q^+\,\zeta}{2\sqrt\xi}\;
\frac{\xi\eta\;\big(\delta^{im}+2i\ee^{im}\lambda_1\big)\,\delta_{\lambda_1\lambda_2}\,t^a_{\alpha\beta}}
{\zeta(1-\zeta)^2\big(\pt^{\,2}+\of^2\big)+\xi\eta\big(\kt^2+\ef^2\big)}
\nonumber\\
&\times\big|\bar q^\beta_{\lambda_2}(\eta q^+,\qq-\kk-\pp)\;g^a_m(\xi q^+,\pp)\;q^\alpha_{\lambda_1}(\zeta q^+,\kk)\big\rangle\,.
\end{align}
The denominator is a single quadratic form in $(\kt,\pt)$, whose constant term is given by
eq.~\eqref{eq:Ef}. Its Fourier transform is therefore the massless one with $Q\to\Ef$:
\begin{equation}
(1-\zeta)\int\frac{\dd^2\kt}{2\pi}\frac{\dd^2\pt}{2\pi}\;
\frac{e^{i\kt\cdot\R}\,e^{-i\pt\cdot\Y}}{\zeta(1-\zeta)^2\big(\pt^{\,2}+\of^2\big)+\xi\eta\big(\kt^2+\ef^2\big)}
=\frac{\Ef\,K_1(\Ef\,\cD)}{\cD}\equiv\cG_E\,,
\label{eq:GE}
\end{equation}
with $\cD$ given by eq.~\eqref{eq:D}. The three partons scatter, so that the Wilson-line structure is the
same as in eq.~\eqref{eq:wf-a}:
\begin{align}
\big|\gamma_T^i(Q,q^+,\w)\big\rangle^{({\rm inst})}_{\qqg}=&\frac{e\,e_f\,g\,q^+}{2(2\pi)^4}\int_{\x,\y,\z}
\int_0^1\dd\zeta\int_0^{1-\zeta}\dd\xi\;\frac{\zeta\,\eta\,\sqrt\xi}{1-\zeta}\;\delta^{(2)}(\w-\cc)
\nonumber\\
&\times\big(\delta^{im}+2i\ee^{im}\lambda_1\big)\,\delta_{\lambda_1\lambda_2}\;\cG_E\;
\Big[U^{ab}(\z)\,V(\x)\,t^b\,V^\dagger(\y)-t^a\Big]_{\alpha\beta}
\nonumber\\
&\times\big|\bar q^\beta_{\lambda_2}(\eta q^+,\y)\;g^a_m(\xi q^+,\z)\;q^\alpha_{\lambda_1}(\zeta q^+,\x)\big\rangle\,.
\label{eq:wf-inst}
\end{align}
There is no instantaneous counterpart of the graph in figure~\ref{fig:qbar}(b), since an instantaneous
propagator cannot cross the shock wave.

\subsubsection{Summary: effective vertices}
\label{sec:eff}

The three contributions \eqref{eq:wf-a}, \eqref{eq:wf-b} and \eqref{eq:wf-inst} can be combined in a
single expression,
\begin{align}
\big|\gamma_T^i(Q,q^+,\w)\big\rangle^{(\bar q)}_{\qqg}=&-\frac{e\,e_f\,g\,q^+}{2(2\pi)^4}\int_{\x,\y,\z}
\int_0^1\dd\zeta\int_0^{1-\zeta}\dd\xi\;\frac{\zeta}{\sqrt\xi}\;\delta^{(2)}(\w-\cc)
\nonumber\\
&\times\bigg\{\cM^{im}_{1\,\lambda_1\lambda_2}\,\Big[U^{ab}(\z)\,V(\x)\,t^b\,V^\dagger(\y)-t^a\Big]_{\alpha\beta}
\nonumber\\
&\qquad\;-\cM^{im}_{2\,\lambda_1\lambda_2}\,\Big[V(\x)\,V^\dagger(\yp)\,t^a-t^a\Big]_{\alpha\beta}\bigg\}
\nonumber\\
&\times\big|\bar q^\beta_{\lambda_2}(\eta q^+,\y)\;g^a_m(\xi q^+,\z)\;q^\alpha_{\lambda_1}(\zeta q^+,\x)\big\rangle\,,
\label{eq:state-qbar}
\end{align}
with the effective vertices
\begin{align}
\cM^{im}_{1\,\lambda_1\lambda_2}(\x,\y,\z)&=\Psi^{im}_{\lambda_1\lambda_2}[\cG](\R,\Y)
-\delta_{\lambda_1\lambda_2}\,\frac{\xi\,\eta}{1-\zeta}\,\big(\delta^{im}+2i\ee^{im}\lambda_1\big)\,\cG_E\,,
\nonumber\\
\cM^{im}_{2\,\lambda_1\lambda_2}(\x,\y,\z)&=\Psi^{im}_{\lambda_1\lambda_2}[\cA](\R,\Y)\,.
\label{eq:M}
\end{align}
The superscript $(\bar q)$ recalls that the gluon is emitted by the antiquark.

A comment is in order about the ultraviolet behaviour at $\Y\to0$, i.e.\ when the gluon is close to its
emitter. Each of the two terms of eq.~\eqref{eq:state-qbar} is singular in this limit, but their
Wilson-line structures become identical and the singularities cancel in the difference, provided that
the radial functions have the same singular behaviour. This is indeed the case: we find
\begin{align}
&\Y^2\,\cG^{(11)}\to\frac{\ef K_1(\ef R)}{\zeta(1-\zeta)R}\,,\qquad
\Y^2\,\cG^{(01)}\to\frac{K_0(\ef R)}{\zeta(1-\zeta)}\,,\qquad
\nonumber\\
&\cG^{(10)}-\cA^{(10)}\to{\rm finite}\,,\qquad\cG^{(00)}-\cA^{(00)}\to{\rm finite}\,,
\qquad\text{for }\Y\to0\,,
\label{eq:UV}
\end{align}
which are the limits of $\Y^2\cA^{(11)}$ and $\Y^2\cA^{(01)}$; the functions $\cG^{(10)}$ and
$\cG^{(00)}$ are only logarithmically singular. The instantaneous term is regular at $\Y\to0$.

\begin{figure}[t]
\centering
\begin{tikzpicture}[scale=0.92]
\begin{scope}
\shock{0.55}
\draw[photon] (-2.3,0)--(-1.3,0);
\draw[quark] (-1.3,0)--(1.7,0.95) node[right]{\small $\x$};
\draw[quark] (-1.3,0)--(1.7,-0.95) node[right]{\small $\y$};
\draw[gluon] (-0.45,0.27)--(1.7,0.15) node[right]{\small $\z$};
\node at (-0.25,0.62) {\small $\xp$};
\node at (-0.3,-1.6) {(a)};
\end{scope}
\begin{scope}[xshift=5.1cm]
\shock{-0.15}
\draw[photon] (-2.3,0)--(-1.3,0);
\draw[quark] (-1.3,0)--(1.7,0.95) node[right]{\small $\x$};
\draw[quark] (-1.3,0)--(1.7,-0.95) node[right]{\small $\y$};
\draw[gluon] (0.55,0.585)--(1.7,0.15) node[right]{\small $\z$};
\node at (-0.75,0.47) {\small $\xp$};
\node at (-0.3,-1.6) {(b)};
\end{scope}
\begin{scope}[xshift=10.2cm]
\shock{0.55}
\draw[photon] (-2.3,0)--(-1.3,0);
\draw[quark] (-1.3,0)--(1.7,-0.95) node[right]{\small $\y$};
\draw[quark] (-1.3,0)--(-1.3,0.5);
\draw[line width=1.6pt] (-1.42,0.25)--(-1.18,0.25);
\draw[quark] (-1.3,0.5)--(1.7,0.95) node[right]{\small $\x$};
\draw[gluon] (-1.3,0.5)--(1.7,0.15) node[right]{\small $\z$};
\node at (-0.3,-1.6) {(c)};
\end{scope}
\end{tikzpicture}
\caption{Gluon emission by the quark: (a) before the collision, (b) after the collision, where the
quark--antiquark pair scatters at the transverse positions $\xp$ and $\y$, and (c) instantaneous
graph.}
\label{fig:q}
\end{figure}

\subsection{Gluon emission by the quark}
\label{sec:q}

The contribution of the graphs in figure~\ref{fig:q} is obtained from eq.~\eqref{eq:state-qbar} by
exchanging the roles of the quark and of the antiquark while keeping the final state fixed: one
exchanges $\zeta\leftrightarrow\eta$, $\x\leftrightarrow\y$, $\lambda_1\leftrightarrow\lambda_2$ and
$\alpha\leftrightarrow\beta$, takes the complex conjugate of the integrand, and changes the overall
sign~\cite{Iancu:2022gpw}. The transverse coordinates transform as
\begin{equation}
\yp\to\xp\equiv\frac{\zeta\x+\xi\z}{\zeta+\xi}\,,\qquad
\R=\x-\yp\;\to\;\y-\xp\equiv-\Rt\,,\qquad\Y=\y-\z\;\to\;\x-\z\equiv\X\,,
\label{eq:coords-q}
\end{equation}
where $\xp$ is the position of the quark before the emission, $\Rt=\xp-\y$ is the size of the pair
produced by the photon and $\X$ the separation between the gluon and the final quark. The
helicity-conserving vertices transform as
\begin{align}
\varphi^{ij}_{\lambda_1\lambda}(\zeta)&\to\varphi^{ij\,*}_{\lambda_2\lambda}(\eta)
=-\varphi^{ij}_{\lambda\lambda_2}(\zeta+\xi)\,,\qquad
\tau^{mn}_{\lambda\lambda_2}(\xi,\eta)\to\tau^{mn\,*}_{\lambda\lambda_1}(\xi,\zeta)
=\tau^{nm}_{\lambda_1\lambda}(\xi,\zeta)\,,
\nonumber\\
\tau^{nm}_{\lambda_1\lambda}(\xi,\zeta)&=\delta_{\lambda_1\lambda}\big[(2\zeta+\xi)\,\delta^{nm}
+2i\,\xi\,\ee^{nm}\lambda\big]\,,
\label{eq:vert-q}
\end{align}
while the helicity-flip vertices become $\mf\,\mu^i_{\lambda\lambda_2}$ at the photon vertex and
$\xi\,\oft\,\mu^m_{\lambda_1\lambda}$ at the gluon vertex. The scales transform as
$\ef\to\eft$, $\of\to\oft$ and $a\to\tilde a$, cf.\ eq.~\eqref{eq:notation-q}, whereas $\cD$, $\Ef$ and
hence $\cG_E$ are invariant. The radial functions become
\begin{equation}
\tG^{(\alpha\beta)}(\Rt,\X)=\frac{\tilde a^{\,\beta}}{2\eta(1-\eta)}\int_0^1\frac{\dd u}{u^{1+\beta}}
\left(\frac{\tilde A_u}{\tilde\Delta_u}\right)^{\frac{\alpha+\beta}{2}}
K_{\alpha+\beta}\Big(\sqrt{\tilde A_u\,\tilde\Delta_u}\Big)\,,
\label{eq:Gtab}
\end{equation}
with
\begin{equation}
\tilde A_u=\eta(1-\eta)Q^2+\mf^2\left(1+\frac{u\,\eta\,\xi}{\zeta}\right)\,,\qquad
\tilde\Delta_u=\Rt^{\,2}+\frac{\tilde a\,\X^2}{u}\,,
\end{equation}
and
\begin{align}
\tA^{(11)}&=\frac{\eft K_1(\eft\tilde R)}{\eta(1-\eta)\,\tilde R}\;\frac{\oft K_1(\oft X)}{X}\,,
&\tA^{(10)}&=\frac{\eft K_1(\eft\tilde R)}{\eta(1-\eta)\,\tilde R}\;K_0(\oft X)\,,
\nonumber\\
\tA^{(01)}&=\frac{K_0(\eft\tilde R)}{\eta(1-\eta)}\;\frac{\oft K_1(\oft X)}{X}\,,
&\tA^{(00)}&=\frac{K_0(\eft\tilde R)}{\eta(1-\eta)}\;K_0(\oft X)\,.
\label{eq:Atexplicit}
\end{align}
Performing these transformations in eq.~\eqref{eq:state-qbar}, we find
\begin{align}
\big|\gamma_T^i(Q,q^+,\w)\big\rangle^{(q)}_{\qqg}=&+\frac{e\,e_f\,g\,q^+}{2(2\pi)^4}\int_{\x,\y,\z}
\int_0^1\dd\zeta\int_0^{1-\zeta}\dd\xi\;\frac{\eta}{\sqrt\xi}\;\delta^{(2)}(\w-\cc)
\nonumber\\
&\times\bigg\{\tM^{im}_{1\,\lambda_1\lambda_2}\,\Big[U^{ab}(\z)\,V(\x)\,t^b\,V^\dagger(\y)-t^a\Big]_{\alpha\beta}
\nonumber\\
&\qquad\;-\tM^{im}_{2\,\lambda_1\lambda_2}\,\Big[t^a\,V(\xp)\,V^\dagger(\y)-t^a\Big]_{\alpha\beta}\bigg\}
\nonumber\\
&\times\big|\bar q^\beta_{\lambda_2}(\eta q^+,\y)\;g^a_m(\xi q^+,\z)\;q^\alpha_{\lambda_1}(\zeta q^+,\x)\big\rangle\,,
\label{eq:state-q}
\end{align}
with
\begin{align}
\tM^{im}_{1\,\lambda_1\lambda_2}(\x,\y,\z)&=\tPsi^{im}_{\lambda_1\lambda_2}[\tG](\Rt,\X)
-\delta_{\lambda_1\lambda_2}\,\frac{\zeta\,\xi}{\zeta+\xi}\,\big(\delta^{im}-2i\ee^{im}\lambda_1\big)\,\cG_E\,,
\nonumber\\
\tM^{im}_{2\,\lambda_1\lambda_2}(\x,\y,\z)&=\tPsi^{im}_{\lambda_1\lambda_2}[\tA](\Rt,\X)\,,
\label{eq:Mt}
\end{align}
and
\begin{align}
\tPsi^{im}_{\lambda_1\lambda_2}[\tF](\Rt,\X)\equiv\sum_\lambda\Big[&
\tau^{nm}_{\lambda_1\lambda}(\xi,\zeta)\,\varphi^{ij}_{\lambda\lambda_2}(\zeta+\xi)\,\tilde R^jX^n\,\tF^{(11)}
+i\,\xi\,\oft\;\mu^m_{\lambda_1\lambda}\,\varphi^{ij}_{\lambda\lambda_2}(\zeta+\xi)\,\tilde R^j\,\tF^{(10)}
\nonumber\\
&-i\,\mf\;\tau^{nm}_{\lambda_1\lambda}(\xi,\zeta)\,\mu^i_{\lambda\lambda_2}\,X^n\,\tF^{(01)}
+\mf\,\xi\,\oft\;\mu^m_{\lambda_1\lambda}\,\mu^i_{\lambda\lambda_2}\,\tF^{(00)}\Big]\,.
\label{eq:Psit}
\end{align}
As in the massless case, the Wilson-line structure of the graphs in which the gluon is emitted before
the collision is the same for the two emitters, whereas for an emission after the collision the colour
matrix now stands to the left of the Wilson lines; the latter structure is again obtained from the
former by letting $\x,\z\to\xp$, because $U^{ab}(\xp)V(\xp)t^b=t^aV(\xp)$. The relative sign and the
relative normalisation of eqs.~\eqref{eq:state-qbar} and \eqref{eq:state-q}, which control the
interference terms in the cross-section, have been checked against the complete Feynman amplitude
(appendix~\ref{app:checks}).

The complete $\qqg$ component of the outgoing state of the transverse photon is the sum
\begin{equation}
\big|\gamma_T^i(Q,q^+,\w)\big\rangle_{\qqg}=\big|\gamma_T^i(Q,q^+,\w)\big\rangle^{(\bar q)}_{\qqg}
+\big|\gamma_T^i(Q,q^+,\w)\big\rangle^{(q)}_{\qqg}\,.
\label{eq:state}
\end{equation}
